% Chapter 1 -- The run-time library (model: TP 6.0 LR Ch 1 [bor s4.3,
% s6.1]): the self-describing Sample procedure, then one AN-12 entry per
% routine, A-Z.  PUB-REF-0001 AN-12, PG-16 guide words, PG-31 letter tabs,
% CP-18 (one entry per shipped RTL routine).
% Sources (Law 1): os/tps/tps.pas ParseWrite/ParseWriteArg (argument
% types, TRUE/FALSE, literal-first rule), ParseFactor Ord/Chr/Length, the
% file-call checks (argument counts and forms; Reset/Rewrite size defaults
% to 1), NameIsFileBuiltin; seed/rt/fileio.asm (Assign closes and copies
% the ASCIIZ name; Reset AH=3Dh read-only; Rewrite AH=3Ch create/truncate;
% BlockRead AH=3Fh; BlockWrite AH=40h; record size 1; negative count;
% FILEIO-ERROR AX= and exit 1); seed/rt/start*.asm serial_* (COM1).  Every
% rule was checked by compiling probes with the shipped compiler; the
% examples are replayed by test-pubs-tps (Free Pascal; the file examples run
% on InitechDOS).
\Gates{test-tps-type-fpc, test-tps-gen-fpc, test-seed-fileio-fpc, test-compiler, test-compiler-os, test-pubs-tps}
\LibraryHeads
\LetterTabs
\chapter{The run-time library}
\ix{run-time library}

This chapter describes every standard procedure and function of Turbo
Initech in alphabetical order.  The first entry is a sample that shows how
the entries are arranged.

\begin{LibEntry}{Sample}{procedure}{}
\LFunction What the routine does, in one sentence.
\LDeclaration The heading of the routine, as it would be declared in
Pascal.  The standard routines are built into the compiler, and some of
them accept arguments that a declared routine could not; the Remarks say
so.
\LResultType For a function, the type of the value it returns.
\LRemarks What each parameter means, what the routine does in every case,
and what happens when it fails.
\LRestrictions What the routine cannot be used for.
\LArrow A remark that you should not overlook.
\LSeeAlso Related routines.
\LExample A complete program that uses the routine, and the output it
produces.  You can compile every example as it stands.
\end{LibEntry}

\begin{LibEntry}{Assign}{procedure}{assign}
\LFunction Names the InitechDOS file that a file variable is to use.
\LDeclaration \Decl{\textbf{procedure} Assign(\textbf{var} F: \textbf{file}; Name: \textbf{string})}
\LRemarks \textit{F} is a variable of type \rw{file}, named by itself.
\textit{Name} is an expression of type \rw{string} or \rw{char}: the file
name, with drive and directory if needed, in the form InitechDOS accepts.
\std{assign}{Assign} records the name for the next \std{reset}{Reset} or
\std{rewrite}{Rewrite} of \textit{F}; it does not open the file, and does not
check that the file exists.  If \textit{F} is open, \std{assign}{Assign}
closes it first.
\LRestrictions \std{assign}{Assign} cannot be used in an expression.
\LSeeAlso \std{reset}{Reset}, \std{rewrite}{Rewrite}
\LExample
\Program{assign}
\ProgramOutput{assign}
\end{LibEntry}

\begin{LibEntry}{BlockRead}{procedure}{blockread}
\LFunction Reads bytes from a file into a string variable.
\LDeclaration \Decl{\textbf{procedure} BlockRead(\textbf{var} F: \textbf{file}; \textbf{var} Buf: \textbf{char}; Count: \textbf{integer}; \textbf{var} Result: \textbf{integer})}
\LRemarks \textit{F} is a file opened by \std{reset}{Reset}.  \textit{Buf} is
a character of a string variable, such as \code{S[1]}: the first byte read
is stored there and the following bytes after it.  \textit{Count} is the
number of bytes to read.  \textit{Result}, an integer variable, receives the
number of bytes actually read, which is less than \textit{Count} when the
end of the file is reached and 0 when the file is already at its end.
The bytes are read from the current position of the file, which then
advances by \textit{Result}.

\std{blockread}{BlockRead} does not change the length of the string that
holds \textit{Buf}.  Set the length first, for example by assigning a string
of the right length, so that the characters read are part of the string's
value.

If InitechDOS reports an error, or \textit{Count} is negative, the program
ends with the message \code{FILEIO-ERROR AX=}\textit{code} on COM1 and exit
code 1.
\LRestrictions \textit{Buf} must be a character of a variable of a
\rw{string} type.  \textit{Result} must be an integer variable named by itself.
\LArrow\std{blockread}{BlockRead} does not check that \textit{Count} bytes fit
in the string after \textit{Buf}; the bytes beyond its capacity overwrite
whatever follows it in memory.
\LSeeAlso \std{blockwrite}{BlockWrite}, \std{reset}{Reset}
\LExample
\Program{blockrd}
\ProgramOutput{blockrd}
The 16 bytes of the file are read in pieces of five; the last call reads
the one byte that remains.
\end{LibEntry}

\begin{LibEntry}{BlockWrite}{procedure}{blockwrite}
\LFunction Writes bytes from a string variable to a file.
\LDeclaration \Decl{\textbf{procedure} BlockWrite(\textbf{var} F: \textbf{file}; \textbf{var} Buf: \textbf{char}; Count: \textbf{integer}; \textbf{var} Result: \textbf{integer})}
\LRemarks \textit{F} is a file opened by \std{rewrite}{Rewrite}.  The
\textit{Count} bytes starting at \textit{Buf}, a character of a string
variable, are written at the current position of the file, which then
advances.  \textit{Result} receives the number of bytes written.

If InitechDOS reports an error, or \textit{Count} is negative, the program
ends with the message \code{FILEIO-ERROR AX=}\textit{code} on COM1 and exit
code 1.
\LRestrictions As for \std{blockread}{BlockRead}.
\LSeeAlso \std{blockread}{BlockRead}, \std{rewrite}{Rewrite}
\LExample
\Program{blockwr}
\ProgramOutput{blockwr}
\end{LibEntry}

\begin{LibEntry}{Chr}{function}{chr}
\LFunction Returns the character with a given code.
\LDeclaration \Decl{\textbf{function} Chr(I: \textbf{integer}): \textbf{char}}
\LResultType \rw{char}
\LRemarks The result is the character whose code is the low eight bits of
\textit{I}, that is, \textit{I} \rw{mod} 256 for a non-negative \textit{I}.
\code{Chr(65)} is \code{'A'}; \code{Chr(321)} is \code{'A'} as well.
\std{chr}{Chr} is a reserved word, written like a function call.
\LSeeAlso \std{ord}{Ord}
\LExample
\Program{chr}
\ProgramOutput{chr}
\end{LibEntry}

\begin{LibEntry}{Length}{function}{length}
\LFunction Returns the current length of a string.
\LDeclaration \Decl{\textbf{function} Length(S: \textbf{string}): \textbf{integer}}
\LResultType \rw{integer}
\LRemarks \textit{S} is an expression of a \rw{string} type.  The result is
the number of characters in its current value, from 0 to the capacity of
the string.  \std{length}{Length} is a reserved word, written like a
function call.
\LRestrictions The argument must be a string; a character, including a
one-character literal such as \code{'A'}, is not accepted.
\LSeeAlso \std{ord}{Ord}
\LExample
\Program{length}
\ProgramOutput{length}
\id{Name} has capacity 10, so the assignment keeps the first ten characters.
\end{LibEntry}

\begin{LibEntry}{Ord}{function}{ord}
\LFunction Returns the ordinal number of a value.
\LDeclaration \Decl{\textbf{function} Ord(X): \textbf{integer}}
\LResultType \rw{integer}
\LRemarks \textit{X} is an expression of type \rw{char}, \rw{boolean} or
\rw{integer}.  For a character the result is its code, 0 through 255; for a
Boolean value it is 0 for \rw{false} and 1 for \rw{true}; an integer is
returned unchanged.  \std{ord}{Ord} is a reserved word, written like a
function call.
\LRestrictions Strings, records, arrays and files are not accepted.
\LSeeAlso \std{chr}{Chr}
\LExample
\Program{ord}
\ProgramOutput{ord}
\end{LibEntry}

\begin{LibEntry}{Reset}{procedure}{reset}
\LFunction Opens an existing file for reading.
\LDeclaration \Decl{\textbf{procedure} Reset(\textbf{var} F: \textbf{file} {[}; Size: \textbf{integer}{]})}
\LRemarks \textit{F} is a file variable to which a name has been given by
\std{assign}{Assign}.  If \textit{F} is open, \std{reset}{Reset} closes it,
then opens the named file for reading at its first byte.  \textit{Size} is
the record size; it may be left out, in which case it is 1, and when it is
given it must have the value 1.

If the file does not exist, or InitechDOS reports another error, the
program ends with the message \code{FILEIO-ERROR AX=}\textit{code} on COM1
and exit code 1; a \textit{Size} other than 1 gives code 1.
\LRestrictions A file opened with \std{reset}{Reset} is opened for reading
only.
\LSeeAlso \std{assign}{Assign}, \std{blockread}{BlockRead},
\std{rewrite}{Rewrite}
\LExample
\Program{reset}
\ProgramOutput{reset}
The second \std{reset}{Reset} moves back to the beginning of the file.

When the file does not exist, the program stops.  On InitechDOS, with no
file \file{LEDGER.OLD} in the current directory, this program

\DosProgram{resetmiss}

prints:

\ProgramOutput{resetmiss}
\end{LibEntry}

\begin{LibEntry}{Rewrite}{procedure}{rewrite}
\LFunction Creates a file, or empties an existing one, and opens it for
writing.
\LDeclaration \Decl{\textbf{procedure} Rewrite(\textbf{var} F: \textbf{file} {[}; Size: \textbf{integer}{]})}
\LRemarks \textit{F} is a file variable to which a name has been given by
\std{assign}{Assign}.  If \textit{F} is open, \std{rewrite}{Rewrite} closes it.
The named file is created; if it exists, its contents are discarded.
\textit{Size} is as for \std{reset}{Reset}.

If InitechDOS reports an error, the program ends with the message
\code{FILEIO-ERROR AX=}\textit{code} on COM1 and exit code 1.
\LArrow Rewriting an existing file destroys its contents.
\LSeeAlso \std{assign}{Assign}, \std{blockwrite}{BlockWrite},
\std{reset}{Reset}
\LExample
\Program{rewrite}
\ProgramOutput{rewrite}
After the second \std{rewrite}{Rewrite} the file holds only the three bytes
written last.
\end{LibEntry}

\begin{LibEntry}{Write}{procedure}{write}
\LFunction Writes values to the serial port COM1.
\LDeclaration \Decl{\textbf{procedure} Write(V1 {[}, V2, ..., Vn{]})}
\LRemarks Each argument is a string literal or an expression of type
\rw{integer}, \rw{boolean}, \rw{char} or \rw{string}.  The arguments are
written one after another without separators: an integer in decimal, with
a minus sign if it is negative; a Boolean value as \code{TRUE} or
\code{FALSE}; a character or a string as it is.  \std{write}{Write} is a
reserved word with a statement form of its own (the \doctitle{Programmer's
Guide}, Chapter 6).  The output goes to COM1, not to the display.
\LRestrictions An argument that begins with a string literal is that
literal alone; put an expression such as \code{'A' + S} in parentheses.
Records, arrays and files cannot be written.  There are no field widths.
\LSeeAlso \std{writeln}{Writeln}
\LExample
\Program{write}
\ProgramOutput{write}
\end{LibEntry}

\begin{LibEntry}{Writeln}{procedure}{writeln}
\LFunction Writes values to the serial port COM1 and ends the line.
\LDeclaration \Decl{\textbf{procedure} Writeln{[}(V1 {[}, V2, ..., Vn{]}){]}}
\LRemarks \std{writeln}{Writeln} writes its arguments as \std{write}{Write}
does, then a line-feed character (code 10).  With no arguments, with or
without the parentheses, it writes only the line feed.
\LRestrictions As for \std{write}{Write}.
\LSeeAlso \std{write}{Write}
\LExample
\Program{writeln}
\ProgramOutput{writeln}
\end{LibEntry}
